% !TeX program = xelatex
% !TeX encoding = UTF-8
% Complete transcription of four handwritten pages; wording and doses retained.
% Source order: pic01.jpg = 20261016_211500919.jpg (handwritten page 1)
%               pic02.jpg = 20261016_211512762.jpg (handwritten page 2)
%               pic03.jpg = 20261016_211518464.jpg (handwritten page 3)
%               pic04.jpg = 20261016_211522965.jpg (handwritten page 4)
% Compile directly with XeLaTeX, or input into a document loading xepersian.
% No clinical corrections or additional medical recommendations have been added.
\ifdefined\settextfont
  \def\PedAcutePhStandalone{0}
\else
  \def\PedAcutePhStandalone{1}
  \documentclass[12pt,a4paper]{article}
  \usepackage[margin=2cm]{geometry}
  \usepackage{amsmath,amssymb,enumitem}
  \usepackage{xepersian}
  \IfFontExistsTF{Amiri}{\settextfont{Amiri}}{\settextfont{FreeSerif}}
  \setlatintextfont{FreeSerif}
\fi
\providecommand{\en}[1]{\lr{#1}}
\newcommand{\paphunclear}[1]{[خوانش نامطمئن: #1]}
\newcommand{\paphpage}[2]{%
  \clearpage\section*{صفحهٔ #1}
  {\footnotesize\en{#2}\par}\medskip}
\ifnum\PedAcutePhStandalone=1 \begin{document}\fi

\paphpage{۱}{pic01.jpg}
\begin{center}
  {\Large\bfseries\en{Acute Pharyngitis}}\par
  \en{Nelson [22] -- 430}
\end{center}

\subsection*{علل فارنژیت}
\begin{itemize}
\item مواجهه با مواد محرک.
\item علل عفونی.
\item بیماری‌های التهابی و خودایمنی.
\end{itemize}
هدف ما افتراق علل مهم فارنژیت مثل \en{GAS} از علل خودمحدودشونده است.

\subsection*{ویروس‌ها}
\begin{itemize}
\item شایع‌ترین‌ها: آنفلوآنزا، پاراآنفلوآنزا، آدنوویروس، کروناویروس، \en{RSV}، \en{EBV}، \ldots
\item اغلب فارنژیت‌های ویروسی بجز \en{IMN} ناشی از \en{EBV}، خفیف هستند.
\item رینوره و سرفه که تدریجاً تشدید می‌شوند.
\item گلودرد هم در \en{COVID} دیده می‌شود هم در \en{MIS-C} (سندرم التهاب چندسیستمی کودکان).
\item وجود ژنژیوواستوماتیت و وزیکول مطرح‌کنندهٔ \en{HSV} است.
  \begin{itemize}
  \item تب بالا، اودینوفاژی.
  \item می‌تواند تا ۱۴ روز طول بکشد.
  \end{itemize}
\item هرپانژین: انواع خاصی از انتروویروس‌ها.
  \begin{itemize}
  \item ضایعات پاپولووزیکولار منتشر یا \en{ulcer} در خلف فارنکس؛ گلودرد شدید، تب.
  \end{itemize}
\item بیماری دست، پا، دهان:
  \begin{itemize}
  \item وزیکول و اولسر در تمام اوروفارنکس، کف دست و پا و بعضاً تنه و اندام‌ها.
  \item کوکساکی‌ویروس، انتروویروس.
  \end{itemize}
\item تب فارنگوکنژنکتیوال:
  \begin{itemize}
  \item آدنوویروس‌ها.
  \item فارنژیت ظرف ۷ روز و کنژنکتیویت ظرف ۱۴ روز رفع می‌شوند.
  \item خیلی آسان در استخر منتقل می‌شود.
  \end{itemize}
\item سرخک (\en{Measles}): اریتم شدید و منتشر فارنکس \en{+ Koplik spot}.
\item \en{IMN}:
  \begin{itemize}
  \item اسپلنومگالی، لنفادنوپاتی، هپاتومگالی.
  \item در نوجوانان با تونسیلیت اگزوداتیو همراه است.
  \end{itemize}
\end{itemize}

\subsection*{یافته‌ها به نفع علل ویروسی}
کنژنکتیویت، کوریزا، سرفه، اسهال، خشونت صدا، استوماتیت، وجود اگزانتم ویروسی.

\subsection*{یافته‌ها به نفع \en{GAS}}
شروع ناگهانی گلودرد، سن ۵--۱۵ سال، تب، سردرد، تهوع و استفراغ، درد شکم،
اگزودای چرکی تونسیلوفارنژیال، زمستان و اوایل بهار، پتشی روی کام، نبود سرفه.

\paphpage{۲}{pic02.jpg}
\subsection*{فوزوباکتریوم نکروفوروم}
\begin{itemize}
\item یک علت مهم فارنژیت در افراد \paphunclear{مسن} نوجوان و بزرگسال.
\item علائم بالینی تا حدود زیادی مشابه \en{GABHS} است.
\item قابل کشت نیست؛ افتراق با \en{PCR} که در دسترس نیست.
\item با سندرم \en{Lemierre} مرتبط است.
  \begin{itemize}
  \item ترومبوفلبیت سپتیک ورید اینترنال ژوگولار: درد و تورم گردن.
  \end{itemize}
\end{itemize}

\subsection*{دیفتری}
یافتهٔ کلیدی \en{Bull neck} (تورم شدید گردن) است \en{+} سودوممبران خاکستری در فارنکس؛ انسداد تنفسی.

\subsection*{تولارمی اوروفارنژیال}
\begin{itemize}
\item مصرف آب، شیر، گوشت پخته‌نشده.
\item گلودرد شدید، تونسیلیت، آدنیت گردنی، زخم‌های دهانی، سودوممبران.
\end{itemize}

\subsection*{سرفهٔ شدید و مزمن پس از فارنژیت}
مایکوپلاسما نومونیه؛ کلامیدوفیلا نومونیه.

\subsection*{\en{GAS}}
\begin{itemize}
\item زیر ۲--۳ سالگی بسیار ناشایع است. سن شایع آن ۵--۱۵ سالگی می‌باشد.
\item در هر زمانی از سال ولی در زمستان و اواخر بهار شایع‌تر است.
\item به صورت بین برادر--خواهرها و همکلاسی‌ها منتشر می‌شود.
\item عامل ۱۵--۳۰ درصد تمام موارد فارنژیت‌ها در کودکان سن مدرسه.
\item پس از کلونیزه شدن فارنکس:
  \begin{itemize}
  \item ناقل بی‌علامت.
  \item عفونت حاد.
  \end{itemize}
\item دورهٔ کمون ۲--۵ روز؛ شروع ناگهانی گلودرد و تب.
\item زبان توت‌فرنگی: در ابتدا یک غشای سفید روی زبان که بعد از چند روز از بین می‌رود:
  \begin{latin}
  \noindent White Strawberry Tongue $\longrightarrow$ Red Strawberry Tongue.
  \end{latin}
\item به طور شایع لنف‌نود انتریور سرویکال متورم و تندر می‌شود.
\item علائم گوارشی به طور شایع وجود دارند؛ در نبود فارنژیت، علائم گوارشی مربوط به \en{GAS} نیست.
\item درد گوش شایع است ولی \en{TM} اغلب \en{N/L} است.
\item باکتری‌های \en{GAS} که اگزوتوکسین تولید می‌کنند:
  \begin{itemize}
  \item \en{Sandpaper macules (Scarlet Fever)}.
  \item از صورت شروع شده و جنرالیزه می‌شوند.
  \item \en{Circumoral pallor}.
  \end{itemize}
\item پتشی دیستال به تورنیکه: علت: شکنندگی مویرگی.
  \begin{latin}
  \noindent Rumpel--Leeds Phenomenon.
  \end{latin}
\item انسان می‌تواند در زندگی چندین بار به \en{GAS} فارنژیت مبتلا شود.
\end{itemize}

\paphpage{۳}{pic03.jpg}
\subsection*{تشخیص}
\begin{itemize}
\item کرایتریای ارزیابی احتمال \en{GAS}: هر کدام ۱ امتیاز:
  \begin{itemize}
  \item سابقهٔ تب \en{$\geq 38$} درجه.
  \item نبود سرفه.
  \item آدنوپاتی انتریور سرویکال تندر.
  \item اگزودا یا تورم تونسیل.
  \item سن ۳--۱۴ سال.
  \item اگر سن \en{$\geq 45$} سال: ۱ امتیاز منفی.
  \end{itemize}
\item امتیاز \en{$\geq 4$} با احتمال ۷۰\% \en{GAS} همراه است.
\item تست آنتی‌بادی ضد استرپتوکوک جایگاهی در فاز حاد ندارد.
\item روش‌های استاندارد: کشت حلق، \en{RADT}، \en{PCR}.
  \begin{itemize}
  \item کشت حلق: گلد استاندارد است.
  \item \en{RADT} ویژگی بالای ۹۵\% دارد؛ اگر مثبت شود دیگر کشت نمی‌گیریم.
  \item حساسیت آن‌ها از کشت کمتر است؛ اگر منفی شوند: کشت.
  \end{itemize}
\item اگر شرح حال و معاینه به نفع ویروس باشد، نباید تست \en{RADT} یا کشت انجام دهیم.
\item \en{PCR} در کمتر از ۱ ساعت از نمونهٔ سواب \en{Throat} با حساسیت و اختصاصیت \en{$\geq 98\%$}.
\item در \en{IMN}: لنفوسیت آتیپیک فراوان، مونونوکلئوز؛ اسلاید آگلوتینیشن تست مثبت (\en{MSAT}).
\end{itemize}

\subsection*{درمان}
\begin{itemize}
\item استامینوفن یا ایبوپروفن: رفع تب و گلودرد.
\item اسپری و لوزنج بی‌حس‌کننده: بنزوکائین، فنول، منتول.
\item \en{CS} سیستمیک: کودکانی که به علت \en{IMN} دچار اختلال راه هوایی هستند.
\item در این حال نمی‌توان \en{CS} را در اکثر موارد فارنژیت در کودکان توصیه کرد.
\item \en{ABT}:
  \begin{itemize}
  \item \en{GGS}: پنی‌سیلین خوراکی.
  \item \en{GABHS}: اریترومایسین خوراکی.
  \end{itemize}
% GGS and the GABHS/erythromycin association are transcribed as written.
\item \en{GAS} ظرف چند روز به طور خودبه‌خود بهبودی پیدا می‌کند ولی \en{ABT} \paphunclear{درمان} را ۱۲--۲۴ ساعت جلو می‌اندازد.
\item هدف از \en{ABT} جلوگیری از تب روماتیسمی حاد است (\en{ARF})؛ باید ظرف ۱۰ روز اول بیماری شروع شود.
\item \en{ABT} از \en{APSGN} جلوگیری نمی‌کند.
\end{itemize}

\subsection*{در چه مواردی حتی بدون نتیجهٔ تست، \en{ABT} را شروع می‌کنیم؟}
\begin{itemize}
\item تشخیص بالینی تب اسکارلت.
\item کودک علامت‌داری که با فرد مبتلا به عفونت قطعی تماس نزدیک در منزل داشته.
\item سابقهٔ \en{ARF} در بیمار و یا اعضای خانواده.
\end{itemize}

\paphpage{۴}{pic04.jpg}
\begin{itemize}
\item طول مدت درمان با پنی‌سیلین خوراکی یا آموکسی‌سیلین ۱۰ روز است.
\item تست \en{follow-up} لازم نیست.
\item در موارد حساسیت به پنی‌سیلین، درمان ۱۰روزه با نسل اول سفالوسپورین:
  سفالکسین، سفادروکسیل.
  \begin{itemize}
  \item اگر ازدیاد حساسیت تیپ \en{I} (آنافیلاکسی) نبوده باشد.
  \end{itemize}
\item درمان ۱۰روزه با اریترومایسین، کلاریترومایسین، کلیندا.
\item درمان ۵روزه با آزیترومایسین.
\item تتراسایکلین، فلوروکینولون‌ها و سولفونامیدها نباید برای درمان \en{GAS} استفاده شوند.
\end{itemize}

\subsection*{درمان ناقل مزمن}
\begin{itemize}
\item عموماً غیرضروری است.
\item هیچ خطری برای بیمار یا اطرافیان ایجاد نمی‌کند.
\item در فارنژیت ویروسی هم تست این بیماران مثبت می‌شود.
\end{itemize}

\subsection*{اندیکاسیون‌ها برای ریشه‌کنی در ناقلان مزمن}
\begin{itemize}
\item اپیدمی \en{ARF} یا \en{APSGN}.
\item سابقهٔ فردی یا خانوادگی \en{ARF}.
\item \paphunclear{وابستگان نزدیک} با فارنژیت‌های مکرر باکتریال.
\item زمانی که قصد تونسیلکتومی داریم.
\item اضطراب بیمار.
\end{itemize}
\begin{itemize}
\item کلیندامایسین خوراکی برای ۱۰ روز:
  \begin{latin}
  \noindent 20 mg/kg/day divided in 3 doses.
  \end{latin}
\item آموکسی‌کلاو برای ۱۰ روز:
  \begin{latin}
  \noindent 40 mg/kg/day up to 2 gr divided in 3 doses.
  \end{latin}
\end{itemize}

\subsection*{فارنژیت مکرر}
\begin{itemize}
\item در سال گذشته \en{$\geq 7$} نوبت؛ در دو سال گذشته هر کدام \en{$\geq 5$} نوبت؛ در سه سال گذشته هر کدام \en{$\geq 3$} نوبت.
% Crossed-out source text after "سال گذشته": بیش از
\item تونسیلکتومی تا ۲ سال تعداد حملات را کاهش می‌دهد.
\item باید حتماً داکیومنته باشند.
\end{itemize}

\begin{itemize}
\item فارنژیت \en{GAS} مکرر: بسیار نامحتمل است که نشان‌دهندهٔ نقص ایمنی باشد.
\item فارنژیت طول‌کشیده (\en{$\geq 1$} هفته) در \en{IMN} و لمیر.
\end{itemize}

\ifnum\PedAcutePhStandalone=1
  \let\PedAcutePhEnd\enddocument
\else
  \let\PedAcutePhEnd\relax
\fi
\PedAcutePhEnd
